\section{Certified Self-Improvement}
\label{sec:method}

\subsection{Setting and patient-level harm}
\label{sec:harm}

A system $f_0$ is deployed on a new domain.
Patients $i$ are drawn i.i.d.\ from the new domain, and labelled patients are split once into an adaptation set $\A$, a certification set $\B$ ($|\B|=n$) and a locked test set $\C$.
From $\A$ alone we build a finite family $\Lam$ of candidate updates; candidate $\lambda$ produces the system $f_\lambda$.

\begin{definition}[Patient harm]
For patient $i$ with ground-truth lesions $G_i$, let $\mathrm{hit}_i(f)$ be the number of lesions in $G_i$ touched by the prediction of $f$, and $\mathrm{Dice}_i(f)$ its Dice score.
Update $\lambda$ \emph{harms} patient $i$, $H_i(\lambda)=1$, if
\begin{equation}
\mathrm{hit}_i(f_\lambda)<\mathrm{hit}_i(f_0)\;\;\lor\;\;\big(G_i\neq\emptyset\,\land\,\mathrm{Dice}_i(f_\lambda)<\mathrm{Dice}_i(f_0)-\varepsilon\big),
\label{eq:harm}
\end{equation}
with $\varepsilon=0.01$. The risk of $\lambda$ is $R(\lambda)=\Pr[H(\lambda)=1]$.
\end{definition}

Three choices are deliberate.
The unit is the \emph{patient}, the unit of clinical consequence, not the voxel or the lesion.
Harm is \emph{relative to $f_0$}: the question at deployment is whether replacing the current system hurts anyone, not whether the new one is perfect.
And harm is \emph{asymmetric}: a large Dice gain in one patient does not offset a lost lesion in another, so an update with excellent mean Dice can have a high risk.
The rollback $\lambda_0\!:=\!f_0$ has $R(\lambda_0)=0$ by definition.

\subsection{Certification by Learn-then-Test}
\label{sec:ltt}

For each $\lambda\in\Lam$ we test $H_\lambda: R(\lambda)>\delta$ on $\B$.
With $k_\lambda=\sum_{i\in\B}H_i(\lambda)$, the exact binomial p-value
\begin{equation}
p_\lambda=\Pr\big[\mathrm{Bin}(n,\delta)\le k_\lambda\big]
\label{eq:pval}
\end{equation}
is super-uniform under $H_\lambda$~\cite{angelopoulos2021ltt}.
Candidates are organised into $m$ \emph{chains}, each ordered from conservative to aggressive (\eg, an action gate lowered step by step); within a chain we use fixed-sequence testing, stopping at the first non-rejection, and each chain is tested at level $\alpha/m$~\cite{bretz2010multiple}.
Let $\hat\Lam\subseteq\Lam$ be the rejected (certified) set.
The deployed update is the certified candidate with the largest measured benefit on $\B$ (subject to a minimum support), or the rollback if $\hat\Lam=\emptyset$.

\begin{theorem}[Single-round guarantee]
If patients in $\B$ and future patients are exchangeable, and $\Lam$ is constructed without $\B$, then the deployed $\hat\lambda$ satisfies
$\Pr\big[R(\hat\lambda)>\delta\big]\le\alpha$.
\label{thm:ltt}
\end{theorem}
\noindent\emph{Proof sketch.} Every $p_\lambda$ is super-uniform under $H_\lambda$; fixed-sequence testing within chains and Bonferroni across them control the family-wise error at $\alpha$; $\hat\lambda$ is either certified or the zero-risk rollback.
Selection by benefit happens only \emph{among} certified candidates, so it cannot inflate the error (\cref{sec:supp_proof}). \hfill$\square$

\para{Power.} With $k=0$ harms, $H_\lambda$ is rejected only if $(1-\delta)^n\le\alpha/m$.
The smallest certifiable tolerance is therefore $\delta_{\min}=1-(\alpha/m)^{1/n}$: for $n=126$ and $m=8$ chains, $\delta_{\min}=0.040$.
The certificate is honest about sample size: a 1\% guarantee needs several hundred certification patients (\cref{fig:power}a).
Because $\delta$ is a free parameter, sweeping it gives a \emph{risk--benefit dial}: for each tolerance, the best certified update and its benefit.

\para{Why a test, and why per patient.}
Comparing mean Dice on $\B$ (``validate and pick the best'') asks whether an update is better \emph{on average}, so a lost lesion in one patient is outweighed by small gains in many; bounding an expected per-patient loss~\cite{angelopoulos2024crc} allows the same trade.
Our risk is the probability of an \emph{event} that no gain elsewhere offsets, and the rule bounds the probability that this risk exceeds $\delta$ (alternatives in \cref{sec:supp_faq}).
\Cref{alg:csi} summarises the procedure.

\begin{algorithm}[t]
\caption{Certified Self-Improvement (one round)}
\label{alg:csi}
\small
\begin{algorithmic}[1]
\Require deployed $f_0$; disjoint $\A$, $\B$ ($|\B|=n$); tolerance $\delta$; level $\alpha$; harm rule (Eq.~\ref{eq:harm})
\State Build candidates $\Lam$ on $\A$ only; order them into $m$ chains, conservative first
\For{each chain $c=1,\dots,m$}
  \For{$\lambda$ in chain $c$, in order}
    \State $k_\lambda\gets\sum_{i\in\B}H_i(\lambda)$ \Comment{harm relative to $f_0$}
    \State $p_\lambda\gets\Pr[\mathrm{Bin}(n,\delta)\le k_\lambda]$
    \If{$p_\lambda\le\alpha/m$} add $\lambda$ to $\hat\Lam$ \Else{} \textbf{break} \EndIf
  \EndFor
\EndFor
\State \Return $\arg\max_{\lambda\in\hat\Lam}\mathrm{benefit}_\B(\lambda)$ if $\hat\Lam\neq\emptyset$, else $f_0$ (rollback)
\end{algorithmic}
\end{algorithm}

\subsection{Two levels of self-improvement}
\label{sec:levels}

\para{Level 1: the refinement agent.}
On top of the segmenter runs an agent that diagnoses and repairs its output (details in \cref{sec:supp_agent}).
Each connected component of the prediction is described by 26 image and model features plus 21 context features (local tumour-to-background ratio, CT texture, anatomical neighbourhood, left--right symmetry, predicted-probability profile).
Gradient-boosted detectors estimate, per component, the probability of eight failure classes (\eg, physiological false positive, under-segmentation); a hand-built policy maps the predicted class to a repair action (removal, regrowth, threshold refinement), applied when the detector's confidence exceeds a gate.
This hand-built agent is the deployed $f_0$ in the new domain.
A candidate update \emph{adds} actions to it but never removes one (an anchored, one-sided trust region), using detectors refitted on $\A$ with isotonic calibration~\cite{guo2017calibration} on cross-fitted predictions, an out-of-distribution fence on the number of components per scan, and optional guards (agreement of old and new detectors; never act on a scan's only component).
The family is
$\Lam_{\text{agent}}=\{\text{detector}\}\times\{\text{guard}\}\times\{\text{gate}\}$, with 2 detectors, 4 guards and 11 gates ($0.99\to0.85$), \ie, $m=8$ chains and 88 candidates.
Per-patient harm on $\B$ is computed from the measured outcome of every action on every component.

\para{Level 2: the segmenter.}
The segmenter is fine-tuned on $\A$ (5 folds; each fold initialised from the corresponding source-domain model, all weights including the head; 100 epochs; nnU-Net defaults otherwise~\cite{isensee2021nnunet}).
The 5-fold ensemble gives the candidate FT, and FT followed by the hand-built agent gives FT+agent; $\Lam_{\text{seg}}=\{\text{FT},\text{FT+agent}\}$, each tested at $\alpha/2$ against the system deployed after Level 1.
Out-of-fold predictions on $\A$ let the Level-1 machinery be re-run on top of FT.

\para{Sequential deployment.}
Each round certifies against the system deployed by the previous round, so the chain of promises reads $f_0\to f_1\to f_2$.
Each round's guarantee holds at $\alpha$; a union bound gives a joint guarantee at $2\alpha$ for both rounds.
